\documentclass[11pt]{article}
\usepackage[margin=1in]{geometry}
\usepackage{amsmath,amssymb,amsthm,mathtools,microtype}
\usepackage[colorlinks=true,linkcolor=blue,citecolor=blue,urlcolor=blue]{hyperref}
\newtheorem{theorem}{Theorem}[section]
\newtheorem{proposition}[theorem]{Proposition}
\newtheorem{lemma}[theorem]{Lemma}
\newtheorem{remark}[theorem]{Remark}
\numberwithin{equation}{section}
\newcommand{\R}{\mathbb R}
\newcommand{\E}{\mathbb E}
\newcommand{\Prb}{\mathbb P}
\newcommand{\dist}{\operatorname{dist}}
\newcommand{\supp}{\operatorname{supp}}
\newcommand{\Bcal}{\mathcal B}
\DeclareMathOperator*{\esssup}{ess\,sup}
\begin{document}

\title{Uniqueness of positive harmonic profiles outside convex cones\newline
for general stable operators}
\author{Kaiwen Jin\thanks{\href{mailto:yui.ui.siki@gmail.com}{yui.ui.siki@gmail.com}}}
\date{October 15, 2026}
\maketitle

\begin{abstract}
We prove uniqueness, up to a positive factor, of positive harmonic functions
outside a proper closed convex cone with nonempty interior for every
nondegenerate symmetric stable operator. The spectral measure may be singular
or supported on finitely many directions. Each profile is assumed continuous,
zero on the cone, and of polynomial growth strictly below the operator's
order. The proof controls the entire growth norm by one interior value.
This follows by normalizing at a maximizing scale and using a lower barrier
whose exponent is smaller than the chosen growth weight. A comparison of
the resulting normalized profiles then yields proportionality.
\end{abstract}

\section{Introduction}\label{sec:introduction}

Positive harmonic functions that vanish on a cone are natural objects in
the study of nonlocal elliptic equations and their blow-up limits. For
symmetric stable operators, their behavior reflects both the cone geometry
and the directions in which jumps are allowed. A finite spectral measure
describes those directions, and it can be concentrated on a few of them.
The associated interior and boundary regularity theory was developed by
Ros-Oton and Serra \cite{RS}. Their half-space Liouville theorem
\cite[Theorem 4.1]{RS} classifies weak solutions of sub-\(2s\) polynomial
growth in a half-space.

For more general cones, proportionality of positive profiles follows from
boundary Harnack inequalities in the pointwise elliptic setting; see
\cite[Proposition 4.4.5]{FR}. Such an approach does not apply to all stable
operators, since even interior Harnack inequalities can fail for singular
spectral measures. A related uniqueness question is stated in
\cite[Open question 4.4, p.~302]{FR}. Ros-Oton and Weidner
\cite[Theorem 4.6 and Corollary 4.7]{RW} establish comparison of positive
cone profiles under an angular \(L^r\) condition, \(r>n/(2s)\), and
directional monotonicity assumptions. Their argument controls nonlocal tails
before comparing functions near the cone.

Here we establish proportionality for all nondegenerate finite spectral
measures, including atomic ones, for profiles with an explicit growth
exponent below the order of the operator. The geometric assumption is that
the closed convex cone is proper and has nonempty interior. These conditions
specify the version of the uniqueness problem proved below; no conclusion
about other growth classes or cones with empty interior is asserted.

Let \(n\ge2\), \(0<s<1\), and let \(\mu\) be a finite even nonnegative
Borel measure on \(S^{n-1}\) such that, for some \(\lambda,\Lambda>0\),
\begin{equation}\label{eq:ellipticity}
 \mu(S^{n-1})\le\Lambda,\qquad
 \inf_{|e|=1}\int_{S^{n-1}}|e\cdot\theta|^{2s}\,\mu(d\theta)
       \ge\lambda.
\end{equation}
The second condition excludes measures supported in a hyperplane. Define
the positive operator
\begin{equation}\label{eq:operator}
 Lu(x)=\frac12\int_{S^{n-1}}\int_0^\infty
 \bigl(2u(x)-u(x+r\theta)-u(x-r\theta)\bigr)
 \frac{dr}{r^{1+2s}}\,\mu(d\theta).
\end{equation}
For continuous solutions the equation is understood in the viscosity sense.
The growth condition in the following theorem makes its tails integrable.

\begin{theorem}\label{thm:main}
Let \(L\) satisfy \eqref{eq:ellipticity}--\eqref{eq:operator}, and let
\(\Sigma\subsetneq\R^n\) be a closed convex cone with nonempty interior.
Suppose \(w_1,w_2\in C(\R^n)\) are zero on \(\Sigma\), strictly positive
on \(\R^n\setminus\Sigma\), and satisfy
\[
 Lw_i=0\quad\hbox{in }\R^n\setminus\Sigma
\]
in the viscosity sense. If, for each \(i=1,2\), there are finite \(C_i\)
and \(\beta_i<2s\) with
\[
 |w_i(x)|\le C_i(1+|x|)^{\beta_i}\qquad(x\in\R^n),
\]
then \(w_1=cw_2\) on \(\R^n\) for some \(c>0\).
\end{theorem}

The functions need not be homogeneous or monotone. The constants used in
the proof concern a fixed operator and cone; uniformity over all such
operators and cones is unnecessary for the proportionality conclusion.

The difficulty is to obtain comparison constants independent of a profile's
initial tail bound. Write \(\alpha=2s\). A ball barrier has some exponent
\(q<\alpha\), and one can choose a growth weight \(p\) above \(q\) and
both actual profile exponents, while retaining \(p<\alpha\).
Normalize a profile at a scale maximizing its weighted supremum. A compactness
argument and supporting balls then give a distance lower bound for that
normalized function. Evaluation at a fixed interior point recovers the
factor \(R^{q-p}\le1\), and hence controls the whole growth norm by that
single value. This controls tails along every direction charged by \(\mu\).

The required compactness includes the cone boundary. We obtain uniform
boundary smallness from the stopped stable process: starting points that
approach the cone have expected exit times tending to zero. With the growth
bound in hand, an annular comparison in the spirit of \cite{RW} works using
positive spectral mass in outward directions. Scaling makes the comparison
global, and an extremal nonnegative difference proves uniqueness.

Section~\ref{sec:analytic} establishes the solution-concept and stopping
facts. Section~\ref{sec:growth} proves boundary compactness and growth
control, and Section~\ref{sec:rigidity} proves proportionality.
Throughout, \(B_R=B_R(0)\) is an open Euclidean ball, suprema are pointwise,
and the generator of the associated stable process is \(-L\).

\section{Analytic facts and stopped harmonic functions}\label{sec:analytic}

Fix the operator and cone in Theorem~\ref{thm:main}, and put
\(\alpha=2s\) and \(D=\R^n\setminus\Sigma\).
For \(0\le p<\alpha\), let
\[
 \Bcal_p=\left\{u:\sup_{x\in\R^n}
          \frac{|u(x)|}{(1+|x|)^p}<\infty\right\},
 \qquad \|u\|_{\Bcal_p}=
          \sup_x\frac{|u(x)|}{(1+|x|)^p}.
\]
The weight is equivalent to \(1+|x|^p\), the convention in
\cite{FR}. A distributional equation \(Lu=f\) means
\(\int uL\varphi=\int f\varphi\) for every
\(\varphi\in C_c^\infty\) supported in the domain. Viscosity inequalities
use a local \(C^2\) touching test, patched to the original function outside
the touching ball. Upper tests give \(L\le f\), and lower tests give
\(L\ge f\).

The symmetric Levy measure associated with our normalization is
\begin{equation}\label{eq:levy-measure}
 \nu(dz)=r^{-1-\alpha}\,dr\,\mu(d\theta),\qquad z=r\theta,\quad r>0.
\end{equation}
This formula denotes a pushforward measure and permits atoms in \(\mu\).
For a smooth function with growth below \(\alpha\),
\begin{equation}\label{eq:compensated}
 Lu(x)=\int\left(u(x)-u(x+z)
       +\mathbf1_{|z|<1}\nabla u(x)\cdot z\right)\nu(dz).
\end{equation}
Symmetry cancels the odd term and gives exactly \eqref{eq:operator}.

Here are the general-measure analytic results we use. Their applicability
can be checked directly. Write \(M_\mu=\mu(S^{n-1})\), and set
\(\delta=(\lambda/(2\Lambda))^{1/\alpha}\).
For every unit \(e\), the contribution of \(|e\cdot\theta|<\delta\)
to its \(\alpha\)-moment is at most \(\lambda/2\). Consequently
\begin{equation}\label{eq:angular-cap}
 \mu\{|e\cdot\theta|\ge\delta\}\ge\lambda/2,
 \qquad \int|e\cdot\theta|^2\,\mu(d\theta)
                  \ge\delta^2\lambda/2.
\end{equation}
Radial integration gives
\begin{align}
 r^\alpha\nu(B_{2r}\setminus B_r)
       &=\frac{M_\mu}{\alpha}(1-2^{-\alpha}),\label{eq:class-upper}\\
 r^{\alpha-2}\int_{B_r}|e\cdot z|^2\nu(dz)
       &\ge\frac{\delta^2\lambda}{2(2-\alpha)}.\label{eq:class-lower}
\end{align}
Thus \(L\) belongs to the general stable class of
\cite[Definitions 2.1.18 and 2.1.21]{FR}. The class constants may differ
from \(\lambda,\Lambda\), which causes no difficulty for a fixed operator.
In particular, that source permits the measure in \eqref{eq:levy-measure}.

\begin{lemma}[General stable analytic results]\label{lem:inputs}
For the fixed operator \(L\), the following statements hold.
\begin{enumerate}
\item If a continuous \(u\in\Bcal_p\) is distributionally harmonic in a ball, its
convolutions with smooth compactly supported mollifiers are smooth and
strongly harmonic in smaller balls. Interior \(C^s\) estimates there are
controlled by \(\|u\|_{\Bcal_p}\).
\item On a ball, continuous exterior data in \(\Bcal_p\) have a unique
continuous distributionally harmonic extension in the same growth class.
\item On a bounded connected domain, a globally nonnegative continuous
distributionally harmonic function in \(\Bcal_p\) is either positive
throughout that domain or identically zero there.
\item There are \(q\in(s,\alpha)\) and \(\eta\in(0,1)\) such that
\begin{equation}\label{eq:ball-barrier}
 \phi(x)=(1-|x|^2)_+^q,\qquad
 L\phi\le-1\quad\hbox{in }B_1\setminus\overline B_{1-\eta}.
\end{equation}
\end{enumerate}
\end{lemma}

The respective sources are \cite[Lemma 2.2.30 and Theorem 2.4.3]{FR},
\cite[Theorem 2.6.12]{FR}, \cite[Theorem 2.4.15(b)]{FR}, and
\cite[Lemma B.1.4]{FR}. For the interior estimate choose the source's
Lebesgue exponent \(2n/\alpha\); the resulting Holder exponent is
\(\alpha-\alpha/2=s\in(0,1)\). This also covers \(s=1/2\).
The Dirichlet existence result is used only on smooth balls with continuous
data; no regularity of the cone boundary is needed for that use.

We will also use stability of the distributional equation under local
uniform convergence with one common \(\Bcal_p\) bound. To see the relevant
tail estimate for a test \(\varphi\in C_c^\infty\), take \(R\) larger
than twice its support radius. A jump from \(|x|>R\) into that support has
length at least \(R/2\). Changing the order of absolutely integrable
integrals gives
\begin{equation}\label{eq:test-tail}
 \int_{|x|>R}(1+|x|)^p|L\varphi(x)|\,dx
 \le C_\varphi\int_{R/2}^\infty(1+r)^p r^{-1-\alpha}\,dr
 \longrightarrow0.
\end{equation}
On bounded sets \(L\varphi\) is bounded by its local second derivatives
and the finite tail mass. This proves both that the distributional pairing
is defined and that the stated stability holds for this fixed \(L\).

\begin{lemma}[Comparison and solution concepts]\label{lem:comparison}
Let \(\Omega\) be bounded and open. Suppose \(a,b\in C(\R^n)\cap\Bcal_p\),
\(La\le0\le Lb\) in \(\Omega\) in the viscosity sense, and
\(a\le b\) on \(\Omega^c\). Then \(a\le b\) on \(\R^n\).
Moreover, a continuous function of growth below \(\alpha\) is viscosity
harmonic in an open set if and only if it is distributionally harmonic there.
\end{lemma}

\begin{proof}
Choose \(0\le\psi\le1\) smooth and compactly supported, with \(\psi=1\)
on a ball containing \(\overline\Omega\). Every jump contribution to
\(L\psi\) in \(\Omega\) is nonnegative, and sufficiently long jumps give
\(L\psi\ge c>0\) there. For a large radial cutoff \(\chi_R\), equal
to one on \(B_R\) and zero outside \(B_{2R}\), the separated tails give
\[
 L(\chi_Ra)\le e_R,\qquad L(\chi_Rb)\ge-e_R,
 \qquad e_R\le C R^{p-\alpha}\longrightarrow0.
\]
These are viscosity inequalities: near a test point the functions are
unchanged, and the changed tail is bounded by
\(C\int_R^\infty(1+r)^p r^{-1-\alpha}\,dr\).
Fix \(t>0\), take \(2e_R<tc\), and put
\(A=\chi_Ra\), \(B=\chi_Rb+t\psi\).
They are ordered outside \(\Omega\), and
\(LA\le e_R<tc-e_R\le LB\).

If \(\sup(A-B)>0\), maximize
\(A(x)-B(y)-|x-y|^2/(2\varepsilon)\) globally. Compact support gives
maximizers, and exterior order and continuity place both in \(\Omega\)
for sufficiently small \(\varepsilon\). At these points,
\[
 A(x)-A(x+z)-B(y)+B(y+z)\ge0\qquad(z\in\R^n).
\]
Use the quadratic upper test for \(A\) and lower test for \(B\).
Their gradients coincide. Far jumps have nonnegative operator difference;
for jumps of length less than \(r\), the difference is bounded below by
\[
 -\varepsilon^{-1}\int_{|z|<r}|z|^2\nu(dz)
 =-\frac{M_\mu}{\varepsilon(2-\alpha)}r^{2-\alpha}.
\]
Let \(r\downarrow0\) at fixed \(\varepsilon\). This contradicts the
strict negative upper bound \(2e_R-tc\) for the same operator difference.
Thus \(A\le B\). On \(\Omega\), the cutoff is one, and \(t\downarrow0\)
proves comparison.

For a continuous distributional solution, the harmonic mollifications in
Lemma~\ref{lem:inputs}(1) converge locally uniformly with a common
sub-\(\alpha\) growth bound. The viscosity tests pass to the limit:
the patched test controls small jumps, local uniform convergence controls
bounded jumps, and polynomial domination controls the tail.
Thus this solution is viscosity harmonic. Conversely, solve the
distributional Dirichlet problem on a ball compactly contained in the given
open set, using the viscosity solution as exterior data. By the preceding
argument the extension is viscosity harmonic. Comparison in both directions
identifies it with the original function. This proves the converse locally
on every such ball, and hence on the open set.
\end{proof}

In particular, linear combinations of our harmonic functions are harmonic.
Multiplication by a smooth cutoff equal to one near the tested domain changes
the equation only by its explicit tail integral. Addition of a smooth
function changes the right-hand side by its classical \(L\)-image.

To control boundary values uniformly over profiles, we use the symmetric
stable process \(X\) with Levy triplet \((0,0,\nu)\), started at zero.
Existence, the generator on \(C_c^\infty\), the jump decomposition, and the
strong Markov property are given in
\cite[Results 7.5, 6.8, 9.7--9.12, and 4.12]{S}.
Our sign is \(\mathcal A=-L\). Indeed its characteristic exponent is
\[
 \Psi(\xi)=\int(1-\cos(\xi\cdot z))\nu(dz)
 =c_\alpha\int|\xi\cdot\theta|^\alpha\mu(d\theta),
 \quad c_\alpha=\int_0^\infty(1-\cos r)r^{-1-\alpha}\,dr\in(0,\infty).
\]
It follows that \(X_t\) has the law \(t^{1/\alpha}X_1\).
Its semigroup is Feller: for \(f\in C_0(\R^n)\), dominated convergence
shows that \(\E f(x+X_t)\) is continuous and vanishes at infinity, while
uniform continuity gives
\(\|P_tf-f\|_\infty\le\E\omega_f(|X_t|)\to0\).
Thus Dynkin's formula \cite[Lemma 12.2]{S} applies to compact smooth tests
and finite-mean stopping times.

\begin{lemma}[Stopped representation]\label{lem:stopped}
Let \(\Omega\subset D\) be bounded and open, and let
\(u\in C(\R^n)\cap\Bcal_b\), \(0\le b<\alpha\), be harmonic in \(D\).
For \(x\in\Omega\) and
\(\tau_\Omega^x=\inf\{t>0:x+X_t\notin\Omega\}\),
\begin{equation}\label{eq:stopped}
 u(x)=\E u(x+X_{\tau_\Omega^x}).
\end{equation}
The stopped positions in the proof have a uniformly bounded moment of every
fixed order strictly between \(b\) and \(\alpha\).
\end{lemma}

\begin{proof}
We give the stopping and integrability details because the spectral measure
can be singular. The filtration is the usual completed right-continuous
natural filtration. If stopping times \(\sigma_j\uparrow\sigma<\infty\) almost surely,
the strong Markov property implies, for \(a,\varepsilon,M>0\),
\begin{align*}
 &\Prb(|X_\sigma-X_{\sigma_j}|>a,\ \sigma\le M)\\
 &\quad\le\Prb(\sigma-\sigma_j>\varepsilon,\ \sigma\le M)
       +\Prb\left(\sup_{0\le t\le\varepsilon}|X_t|>a\right).
\end{align*}
Let \(j\to\infty\), then \(\varepsilon\downarrow0\), then \(M\to\infty\).
This proves convergence in probability. Cadlag paths already give a left
limit on paths where the sequence strictly announces \(\sigma\), so that
limit equals \(X_\sigma\) almost surely. Where the times eventually equal
\(\sigma\), convergence is immediate. We have proved the required
quasi-left continuity.

For completeness, exits from our bounded open sets are stopping times and
their positions belong to the complement. If \(\Omega\subset B_A\), put
\(F=\Omega^c\) and
\(\sigma_m=\inf\{t>0:\dist(x+X_t,F)<1/m\}\), starting with large \(m\).
Hitting this open neighborhood is a stopping time by rational approximation
and right continuity. Its position has distance at most \(1/m\) from \(F\).
Every \(\sigma_m\) is bounded by \(T_{2A}\), the first jump longer than
\(2A\): a path still in \(B_A\) cannot remain there after such a jump.
This time is exponential, of rate \(M_\mu(2A)^{-\alpha}/\alpha>0\).
The increasing limit \(\sigma\) is finite; quasi-left continuity gives
\(x+X_\sigma\in F\). Since \(\sigma_m\le\tau_\Omega^x\le\sigma\),
the limit is the exit time, with the asserted properties.

Choose \(p'\in(b,\alpha)\) and a nonnegative smooth
\(G\asymp(1+|y|)^{p'}\). For compact smooth cutoffs
\(G_N=\chi_NG\), \(0\le\chi_N\le1\), \(\chi_N=1\) on \(B_N\),
the generators \(\mathcal A G_N\) are uniformly bounded on \(B_A\), for \(N>4A\).
Small jumps are bounded by local second derivatives times \(|z|^2\);
large jumps by \(C(1+|z|)^{p'}\), which is integrable against the tail
of \(\nu\). Dynkin and Fatou therefore give, for every stopping time
\(\sigma\le\tau_\Omega^x\) and finite horizon \(T\),
\begin{equation}\label{eq:exit-moment}
 \E G(x+X_{\sigma\wedge T})
       \le G(x)+C\E T_{2A}<\infty.
\end{equation}
The same estimate holds for exits from inner subsets of \(\Omega\).

Exhaust \(\Omega\) by
\(\Omega_j=\{y\in\Omega:\dist(y,\Omega^c)>1/j\}\), and write
\(\tau_j\) for their exits. Mollify \(u\) at a scale below the distance
of \(\overline\Omega_j\) from \(D^c\). The mollification is smooth and
harmonic on \(\Omega_j\). Apply Dynkin first to its compact smooth cutoffs
and \(\tau_j\wedge T\). The cutoff generator error tends uniformly to
zero, with bound \(C N^{b-\alpha}\).
The moment estimate permits removal of the cutoff and the mollifier.
If \(b=0\), the function values are bounded. If \(b>0\), put
\(r=p'/b>1\); then
\[
 \sup_{j,T}\E|u(x+X_{\tau_j\wedge T})|^r<\infty,
 \qquad
 \E(|Y|\mathbf1_{|Y|>H})\le H^{1-r}\E|Y|^r.
\]
This proves uniform integrability, including for the mollifications.
Consequently \(T\to\infty\) gives
\(u(x)=\E u(x+X_{\tau_j})\).

Finally \(\tau_j\uparrow\tau_\Omega^x\). Otherwise the limiting position,
given by quasi-left continuity, would be interior to \(\Omega\), whereas
each exited inner position has distance at most \(1/j\) from \(\Omega^c\).
Stopped positions converge almost surely; continuity of \(u\) and the
same uniform integrability give \eqref{eq:stopped}.
\end{proof}

\section{Boundary compactness and growth control}\label{sec:growth}

The operator and cone stay fixed throughout this section. Constants may
depend on them, on a chosen \(p<\alpha\), and on the indicated compact sets.
They will be uniform over the functions in each stated family.

\begin{proposition}[Uniform boundary smallness]\label{prop:boundary}
Fix \(0<p<\alpha\) and \(K<\infty\). Consider all nonnegative continuous
functions, zero on \(\Sigma\), harmonic in \(D\), and bounded by
\begin{equation}\label{eq:controlled}
 u(x)\le K(1+|x|)^p.
\end{equation}
For every \(0<A<\infty\), these functions tend uniformly to zero as points
of \(D\cap\overline B_A\) approach \(\Sigma\).
\end{proposition}

\begin{proof}
We first show that \(X_1\) has full support. The support of \(\mu\) spans
\(\R^n\) by directional nondegeneracy. Choose independent support
directions \(\theta_1,\ldots,\theta_n\). Write the center of any prescribed
ball as \(\sum b_j\theta_j\). For nonzero coefficients, small disjoint
angular and radial neighborhoods of \(b_j\theta_j\) have positive
\(\nu\)-mass; use evenness if \(b_j<0\). Requiring one jump in each bin
and no other jump longer than \(\varepsilon\) has positive probability.
Choose the bins so their sum lies in half the prescribed ball, and then
take \(\varepsilon\) below their radial supports. The remaining independent
symmetric small-jump martingale has
\[
 \E|X_1^{\le\varepsilon}|^2
 =\int_{|z|\le\varepsilon}|z|^2\nu(dz)
 =\frac{M_\mu}{2-\alpha}\varepsilon^{2-\alpha}\longrightarrow0.
\]
Chebyshev's inequality gives positive probability that this remaining part
lies in the other half-ball. The jump assertions follow from
\cite[Theorems 9.7, 9.9, 9.12 and Corollary 9.11]{S}.
Symmetry makes the intermediate compensating drift zero. This proves full
support, including for atomic spectral measures.

Since \(\operatorname{int}\Sigma\) is nonempty and invariant under positive
dilations, scaling gives
\[
 \Prb(X_t\in\operatorname{int}\Sigma)
       =\Prb(X_1\in\operatorname{int}\Sigma)=p_0>0\quad(t>0).
\]
The event of visits to this open cone at arbitrarily small positive times
therefore has probability at least \(p_0\). It is a germ event, and has
probability one. Indeed a germ event is independent of the increments after
every \(\varepsilon>0\). As \(\varepsilon\downarrow0\), right continuity
shows independence from every finite collection of path values, hence from
the whole path sigma algebra and from itself.

Let \(\Omega=D\cap B_{2A+2}\), and couple all starting points by the same
process. If \(x_k\in D\cap\overline B_A\) tends to \(y\in\Sigma\), then
\begin{equation}\label{eq:exit-small}
 \tau_\Omega^{x_k}\longrightarrow0\quad\hbox{almost surely},
 \qquad \E\tau_\Omega^{x_k}\longrightarrow0.
\end{equation}
To prove the first assertion, use any of the arbitrarily early times with
\(X_t\in\operatorname{int}\Sigma\). Convexity of the cone gives
\(y+\operatorname{int}\Sigma\subset\operatorname{int}\Sigma\), so
\(x_k+X_t\in\Sigma\) for all sufficiently large \(k\). The second
assertion follows by domination with the same \(T_{4A+4}\).

Choose a smooth nonnegative \(F_A\), zero on \(B_A\), dominating
\(K(1+|x|)^p\) outside \(B_{2A+2}\), with this polynomial growth.
Its generator is bounded on \(B_{2A+2}\). At exit, \(u=0\) on \(\Sigma\)
and \(u\le F_A\) outside the ball. Lemma~\ref{lem:stopped} gives
\begin{equation}\label{eq:boundary-modulus}
 0\le u(x)\le\E F_A(x+X_{\tau_\Omega^x})
       \le C_A\E\tau_\Omega^x,\qquad x\in D\cap B_A.
\end{equation}
For the last inequality apply Dynkin to nonnegative compact cutoffs of
\(F_A\), whose generators have one common bound on the stopped ball,
and then use Fatou. Equations \eqref{eq:exit-small}--\eqref{eq:boundary-modulus}
and compactness of \(\Sigma\cap\overline B_A\) prove the assertion.
\end{proof}

\begin{lemma}[Compactness and positivity]\label{lem:compact}
The family in Proposition~\ref{prop:boundary} is precompact for local uniform
convergence on \(\R^n\). A subfamily normalized by \(\sup_{B_1}u=1\), or
by \(u(x_*)=1\) at a fixed \(x_*\in D\), has a common positive lower bound
on every compact subset of \(D\).
\end{lemma}

\begin{proof}
Interior Holder estimates from Lemma~\ref{lem:inputs}, together with the
uniform smallness near \(\Sigma\), give equicontinuity on every closed ball.
Arzela--Ascoli and a diagonal extraction yield local uniform compactness.
The common growth bound and \eqref{eq:test-tail} preserve distributional
harmonicity. Limits are nonnegative, continuous, and zero on \(\Sigma\).

The set \(D\) is path connected. Choose \(e\in\operatorname{int}\Sigma\).
For \(x\notin\Sigma\), the ray \(x-te\), \(t\ge0\), stays outside
\(\Sigma\), since otherwise adding \(te\in\Sigma\) would put \(x\) in
\(\Sigma\). There is a nonzero functional \(\ell\ge0\) on \(\Sigma\):
project any point outside the cone onto it and take the negative outward
normal. Since \(e\) is interior, \(\ell(e)>0\). For two given points,
their translates by a sufficiently large common \(-te\) lie in the
half-space \(\{\ell<0\}\). Join them there and along their rays.

A nonzero nonnegative limit is consequently positive throughout \(D\),
by the strong maximum principle in Lemma~\ref{lem:inputs}(3), applied on
bounded connected neighborhoods of such paths. Both stated normalizations
pass to local uniform limits and exclude the zero function. If a compact
lower bound failed, a convergent sequence of points on that compact set
and a locally uniformly convergent sequence of profiles would give a
nonzero limit with an interior zero, a contradiction.
\end{proof}

We next turn the controlled polynomial bound into one independent of the
profile's initial growth constant. Let \(q,\eta\) be as in
\eqref{eq:ball-barrier}.

\begin{lemma}[Distance lower bound]\label{lem:distance}
Fix \(0<p<\alpha\). Every nonnegative continuous function, zero on \(\Sigma\)
and harmonic in \(D\), with
\begin{equation}\label{eq:unit-family}
 \sup_{B_1}v=1,\qquad v(x)\le\max(1,|x|^p),
\end{equation}
satisfies
\begin{equation}\label{eq:distance-lower}
 v(x)\ge a\,\dist(x,\Sigma)^q\quad(x\in D\cap B_1),
\end{equation}
for some \(a>0\) independent of \(v\).
\end{lemma}

\begin{proof}
Lemma~\ref{lem:compact} gives a common positive lower bound on
\(K_\eta=\{z\in\overline B_3:\dist(z,\Sigma)\ge\eta\}\).
For \(x\in D\cap B_1\), let \(y\) be its nearest point on \(\Sigma\),
\(d=|x-y|\), and \(\zeta=(x-y)/d\). Convex projection gives
\((x-y)\cdot(z-y)\le0\) for \(z\in\Sigma\).
Taking \(z=0,2y\) yields \((x-y)\cdot y=0\). In particular
\(|y|\le1\), \(0<d\le1\), and
\(\Sigma\subset\{z:\zeta\cdot z\le0\}\).
Thus \(B_1(y+\zeta)\subset D\), and its inner ball of radius
\(1-\eta\) lies in \(K_\eta\).

A sufficiently small fixed multiple \(a\phi(\,\cdot-y-\zeta)\) is below
\(v\) on that inner ball and outside the unit ball. Comparison on the
outer collar, using \eqref{eq:ball-barrier}, makes the same order hold
everywhere. Evaluation at \(x\) gives
\[
 v(x)\ge a(1-|x-y-\zeta|^2)^q
        =a(2d-d^2)^q\ge a d^q.
\]
\end{proof}

For a nonnegative function define its polynomial growth norm by
\begin{equation}\label{eq:growth-norm}
 \|u\|_{p,*}=\sup_{R\ge1}R^{-p}U_u(R),
 \qquad U_u(R)=\sup_{B_R}u.
\end{equation}
This measures the whole tail, including the directions charged by \(\mu\).

\begin{proposition}[Growth controlled by one value]\label{prop:growth}
Fix \(q<p<\alpha\) and \(x_*\in D\cap B_{1/2}\).
For every positive continuous harmonic profile in \(D\), zero on \(\Sigma\),
whose actual polynomial growth exponent is strictly below \(p\),
\begin{equation}\label{eq:growth-control}
 \|u\|_{p,*}\le C u(x_*).
\end{equation}
The constant is independent of the profile and its initial growth constant.
\end{proposition}

\begin{proof}
Write \(U=U_u\). Continuity makes \(U(R)\) continuous, and the strict
growth margin makes \(U(R)/R^p\to0\). Its positive supremum is therefore
attained at some \(R_0\ge1\). The function
\[
 v(x)=\frac{u(R_0x)}{U(R_0)}
\]
satisfies \eqref{eq:unit-family}: the maximizing property gives its bound
for \(|x|\ge1\), while \(U(R_0)\) gives the bound inside the unit ball.
The stable scaling identity
\(L(u(R\,\cdot))(x)=R^\alpha(Lu)(Rx)\) and the cone property preserve
harmonicity and the zero set.

Apply Lemma~\ref{lem:distance} at \(x_*/R_0\). Since the cone distance
scales linearly,
\[
 u(x_*)\ge aU(R_0)R_0^{-q}\dist(x_*,\Sigma)^q.
\]
Consequently
\[
 \|u\|_{p,*}=\frac{U(R_0)}{R_0^p}
 \le a^{-1}\dist(x_*,\Sigma)^{-q}R_0^{q-p}u(x_*)
 \le C u(x_*),
\]
because \(p>q\) and \(R_0\ge1\).
\end{proof}

\section{Comparison and proportionality}\label{sec:rigidity}

Let \(\beta_1,\beta_2\) be the growth exponents of the original profiles
\(w_1,w_2\) in Theorem~\ref{thm:main}. Choose once and for all
\begin{equation}\label{eq:p-choice}
 \max(q,\beta_1,\beta_2)<p<\alpha,
 \qquad x_*\in D\cap B_{1/2}.
\end{equation}
All profiles considered below have actual growth exponent below this same
\(p\). Normalize two of them by \(u(x_*)=v(x_*)=1\).
Proposition~\ref{prop:growth} and Lemma~\ref{lem:compact} give
\begin{equation}\label{eq:normalized-growth}
 u(x)+v(x)\le C(1+|x|)^p,
 \qquad \inf_K v\ge c_K>0\quad(K\Subset D),
\end{equation}
with constants uniform over such normalized profiles.

The next comparison uses a bump outside the ball of interest. Its forcing
comes from an angular cap of positive measure.

\begin{lemma}[Annular forcing]\label{lem:annulus}
There is \(\eta_0\in C_c^\infty(D)\), \(0\le\eta_0\le1\), supported
in an annulus about radius \(3/4\), outside \(B_{2/3}\), such that
\begin{equation}\label{eq:annular-forcing}
 L\eta_0\le-c_0<0\quad\hbox{in }D\cap B_{2/3}.
\end{equation}
\end{lemma}

\begin{proof}
For \(z\in D\cap B_{2/3}\), let \(\zeta\) be the outward unit normal
from its projection onto \(\Sigma\). The projection identities in the
proof of Lemma~\ref{lem:distance} give
\(\Sigma\subset\{\zeta\cdot x\le0\}\) and
\(\zeta\cdot z=\dist(z,\Sigma)>0\).
By \eqref{eq:angular-cap} and evenness,
\[
 \mu\{\theta:\zeta\cdot\theta\ge\delta\}\ge m,
 \qquad \delta=(\lambda/(2\Lambda))^{1/\alpha},\quad m=\lambda/4.
\]
Choose \(2/3<a_-<a_+<4/5\), both near \(3/4\). Every ray
\(z+r\theta\) with \(\zeta\cdot\theta\ge\delta\) crosses this radial
band over an interval with
\[
 a_--2/3\le r\le a_++2/3,\qquad r_+-r_-\ge a_+-a_-.
\]
The length bound follows from the 1-Lipschitz property of the radial norm.
Throughout the crossing,
\(\dist(z+r\theta,\Sigma)\ge\zeta\cdot(z+r\theta)\ge r\delta\).
Take \(0<\rho<\delta(a_--2/3)\), and choose a smooth bump equal to one
on the compact set
\(\{a_-\le|x|\le a_+,\ \dist(x,\Sigma)\ge\rho\}\),
supported in a slightly wider annulus outside \(B_{2/3}\) and in
\(\{\dist(x,\Sigma)>\rho/2\}\).
At every \(z\) in the tested ball its value is zero, and hence
\[
 -L\eta_0(z)=\int\eta_0(z+h)\nu(dh)
 \ge m(a_+-a_-)(a_++2/3)^{-1-\alpha}>0.
\]
\end{proof}

\begin{proposition}[Comparison of normalized profiles]\label{prop:local}
For the fixed \(p,x_*\) in \eqref{eq:p-choice}, there is \(C_1\ge1\)
such that every normalized pair as above satisfies
\begin{equation}\label{eq:local-comparison}
 C_1^{-1}u\le v\le C_1u\quad\hbox{on }B_{1/2}.
\end{equation}
\end{proposition}

\begin{proof}
Take smooth cutoffs \(0\le\chi,\xi\le1\), with \(\chi=1\) on
\(B_{3/4}\), \(\chi=0\) outside \(B_{4/5}\), \(\xi=1\) on
\(B_{1/2}\), and \(\xi=0\) outside \(B_{2/3}\). Set
\[
 f=\chi u+A(\xi-1)+B\eta_0.
\]
In \(D\cap B_{2/3}\), the cutoff identity in the viscosity sense gives
\[
 L(\chi u)(z)=\int(1-\chi(z+h))u(z+h)\nu(dh)\le C.
\]
The uniform bound follows from \eqref{eq:normalized-growth}, the positive
separation of \(B_{2/3}\) from the support of \(1-\chi\), and
\(p<\alpha\). Also \(|L\xi|\le C\).
First take \(A>\sup_{B_1}u\), then choose \(B\) large enough that
\eqref{eq:annular-forcing} gives \(Lf\le-1\) in the tested domain.
Both constants can be chosen uniformly over the normalized family.

For \(|z|\ge2/3\), \(\chi u-A\le0\), so \(f\le B\eta_0\).
On \(\supp\eta_0\Subset D\), \eqref{eq:normalized-growth} supplies a
common positive lower bound for \(v\). Thus \(f\le C_1v\) outside the
tested ball. On \(\Sigma\), \(f\le0=C_1v\).
Lemma~\ref{lem:comparison} on \(D\cap B_{2/3}\) gives \(f\le C_1v\)
everywhere. In \(B_{1/2}\), \(f=u\). Interchanging \(u,v\) proves the claim.
\end{proof}

This comparison concerns entire cone profiles with controlled polynomial
growth. The control at infinity in Proposition~\ref{prop:growth} is what
makes its constant uniform over normalized profiles.

\begin{proof}[Proof of Theorem~\ref{thm:main}]
For any two positive profiles \(u,v\) and \(R\ge1\), apply
Proposition~\ref{prop:local} to
\[
 u_R(x)=\frac{u(Rx)}{u(Rx_*)},\qquad
 v_R(x)=\frac{v(Rx)}{v(Rx_*)}.
\]
These remain in the same class. Evaluation at \(x_*/R\) bounds
\(u(Rx_*)/v(Rx_*)\) between \(C_1^{-1}\) and \(C_1\) times
\(u(x_*)/v(x_*)\). For any fixed \(x\in D\), choose \(R>2|x|\) and
evaluate once more at \(x/R\). We obtain the global comparison
\begin{equation}\label{eq:global-comparison}
 C_1^{-2}\frac{u(x_*)}{v(x_*)}
 \le\frac{u(x)}{v(x)}
 \le C_1^2\frac{u(x_*)}{v(x_*)}\qquad(x\in D).
\end{equation}

Let \(\kappa=\inf_D u/v>0\), and put \(h=u-\kappa v\).
By Lemma~\ref{lem:comparison}, \(h\) is harmonic in \(D\).
It is continuous, nonnegative, zero on \(\Sigma\), and has growth exponent
at most \(\max(\beta_1,\beta_2)<p\). If it is nonzero, the strong maximum
principle and connectedness proved in Lemma~\ref{lem:compact} make it
positive throughout \(D\). Apply \eqref{eq:global-comparison} to \(h,v\).
This forces \(\inf_D h/v>0\), contradicting
\(\inf_D h/v=\inf_D(u/v-\kappa)=0\). Hence \(h=0\), and
\(u=\kappa v\) everywhere. The choices in \eqref{eq:p-choice} are always
possible because each original exponent is strictly less than \(\alpha\).
They impose no further condition on the theorem.
\end{proof}

\paragraph{AI use.}
OpenAI Codex was used to develop the argument, check cited source statements,
draft the manuscript, and perform internal mathematical and editorial review.
The manuscript is submitted for independent mathematical scrutiny.

\begin{thebibliography}{9}
\bibitem{FR}
X. Fernandez-Real and X. Ros-Oton,
\emph{Integro-Differential Elliptic Equations},
manuscript, arXiv:2411.12455v1, 2024.
\href{https://arxiv.org/abs/2411.12455v1}{arxiv.org/abs/2411.12455v1}.

\bibitem{RS}
X. Ros-Oton and J. Serra,
\emph{Regularity theory for general stable operators},
arXiv:1412.3892v1, 2014.
\href{https://arxiv.org/abs/1412.3892v1}{arxiv.org/abs/1412.3892v1}.

\bibitem{RW}
X. Ros-Oton and M. Weidner,
\emph{Obstacle problems for nonlocal operators with singular kernels},
arXiv:2308.01695v1, 2023.
\href{https://arxiv.org/abs/2308.01695v1}{arxiv.org/abs/2308.01695v1}.

\bibitem{S}
R. L. Schilling,
\emph{An Introduction to L\'evy and Feller Processes:
Advanced Courses in Mathematics -- CRM Barcelona 2014},
arXiv:1603.00251v2, 2016.
\href{https://arxiv.org/abs/1603.00251v2}{arxiv.org/abs/1603.00251v2}.
\end{thebibliography}
\end{document}
